% Copyright 2020 by Robert Hildebrand
%This work is licensed under a
%Creative Commons Attribution-ShareAlike 4.0 International License (CC BY-SA 4.0)
%See http://creativecommons.org/licenses/by-sa/4.0/
\chapter{Advanced Models}
In this chapter, we consider some models where binary varibales are used as an "on/off" switch.  








\section{Advanced Models}
\subsection{Work Scheduling}
\subsection{Multi period Models}
\todo[inline]{Fill in this subsection}

\subsubsection{Production Planning}

\subsubsection{Crop Planning}


\subsection{Mixing Problems}

\subsection{Financial Planning}

\todo[inline]{Fill in this subsection}




\subsection{Network Flow}
\begin{resource}
\begin{itemize}
\item \href{https://ocw.mak.ac.ug/courses/electrical-engineering-and-computer-science/6-046j-design-and-analysis-of-algorithms-spring-2012/lecture-notes/MIT6_046JS12_lec13.pdf}{MIT - CC BY NC SA 4.0 license}
\item \href{https://www.cs.princeton.edu/~wayne/kleinberg-tardos/pearson/}{Slides for Algorithms book by Kleinberg-Tardos}
\end{itemize}
\end{resource}





\paragraph{Sets}
A finite network $G$ is described by a finite set of vertices $V$ and a finite set $\mathcal{A}$ of arcs. Each arc $(i,j)$ has two key attributes, namely its tail $j \in V$ and its head $i \in V$. 


We think of a (single) commodity as being allowed to "flow" along each arc, from its tail to its head. 
\paragraph{Variables}
Indeed, we have "flow" variables
$$
x_{ij}:=\text { amount of flow on } \operatorname{arc} (i,j) \text{ from vertex $i$ to vertex $j$},
$$
for all $(i,j) \in \mathcal{A}$. 



\subsubsection{Maximum Flow Problem}

\begin{tcolorbox}
\begin{align}
\max \ \ & \sum_{(s, i) \in \mathcal{A}} x_{si} &
\text{max total flow from source}\\
s.t.\ \ 
&\sum_{i : (i,v) \in \mathcal A} x_{i v}-\sum_{j : (v,j) \in \mathcal{A}} x_{v j} =0  &  v \in V \backslash\{s, t\}\\
%[a constraint per edge and a constraint per nonterminal node]\\
&0 \leq x_{ij} \leq  u_{ij} \quad \forall(i,j) \in \mathcal{A}
\end{align}

\end{tcolorbox}
%\documentclass[border=4mm]{standalone}
%\usepackage{tikz}
%\usetikzlibrary{arrows.meta,positioning}
%\begin{document}

\paragraph{Shortest Path Problem}
~\\

\begin{tcolorbox}
$$
\begin{aligned}
& \text { minimize } & \sum_{u \rightarrow v} \ell_{u \rightarrow v} \cdot x_{u \rightarrow v} & \\
\text { subject to } & \sum_{u} x_{u \rightarrow s}-\sum_{w} x_{s \rightarrow w} &=1 & \\
\sum_{u} x_{u \rightarrow t}-\sum_{w} x_{t \rightarrow w} &=-1 & \\
\sum_{u} x_{u \rightarrow v}-\sum_{w} x_{v \rightarrow w} &=0 & \text { for every vertex } v \neq s, t \\
x_{u \rightarrow v} & \geq 0 & \text { for every edge } u \rightarrow v
\end{aligned}
$$
\end{tcolorbox}

Or maybe write it like this:
\begin{align}
\text{min} \sum_{(i, j) \in A} c_{i j} x_{i j} \\
s.t. 
\sum_{(i, j) \in \delta^{+}(i)} x_{i j} = 0 \forall i \in V\setminus\{s,t\} \\
\sum_{(i, j) \in \delta^{+}(s)} x_{i j} = -1 \\
\sum_{(i, j) \in \delta^{+}(t)} x_{i j} = 1 \\
% _decvars_
x_{i j} \in\{0,1\} \forall(i, j) \in A, 
\end{align}

\subsubsection{Minimum Cost Network Flow}
\paragraph{Parameters}

We assume that flow on arc $(i,j)$ should be non-negative and should not exceed
$$
u_{ij}:=\text { the flow upper bound on } \operatorname{arc} (i,j) \text {, }
$$
for $(i,j) \in \mathcal{A}$. Associated with each arc $(i,j)$ is a cost
$$
c_{ij}:=\text { cost per-unit-flow on arc } (i,j),
$$
for $(i,j) \in \mathcal{A}$. The (total) cost of the flow $x$ is defined to be
$$
\sum_{(i,j) \in \mathcal{A}} c_{ij} x_{ij} .
$$

We assume that we have further data for the nodes. Namely,
$$
b_{v}:=\text { the net supply at node } v,
$$
for $v \in V$. 

A flow is conservative if the net flow out of node $v$, minus the net flow into node $v$, is equal to the net supply at node $v$, for all nodes $v \in V$.

The (single-commodity min-cost) network-flow problem is to find a minimumcost conservative flow that is non-negative and respects the flow upper bounds on the arcs. 





\paragraph{Objective and Constraints}
We can formulate this as follows:
\begin{tcolorbox}
$$
\begin{aligned}
\min & \sum_{(i,j) \in \mathcal{A}} c_{ij}\, x_{ij} & \text{ minimize cost}\\
& \sum_{(i,v) \in \mathcal{A}} x_{iv}-\sum_{(v,i)  \in \mathcal{A} } x_{vi}=b_{v}, \quad \text{ for all }  v \in V,  & \text{ flow conservation}\\
& 0 \leq x_{ij} \leq u_{ij}, \quad \text{ for all }  (i,j) \in \mathcal{A} .
\end{aligned}
$$
\end{tcolorbox}

\begin{theorem}{Integrality of Network Flow}{integralitynetworkflow}
If the capacities and demands are all integer values, then there always exists an optimal solution to the LP that has integer values.
\end{theorem}

\subsection{Multi-Commodity Network Flow}
In the same vein as the Network Flow Problem
\begin{align*}
\min~& \sum_{k=1}^{K} \sum_{e \in \mathcal{A}} c^k_e x^k_e \\
&\sum_{e\in \mathcal{A} ~:~ t(e)=v} x^k_e - \sum_{e\in \mathcal{A} ~:~ h(e)=v} x^k_e 
= b^k_v,~ \mbox{ for } v \in \mathcal{N},~ k=1,2,\ldots,K;\\
& \sum_{k=0}^{K}  x^k_{e} \leq u_e,~ \mbox{ for } e \in \mathcal{A};\\
&  x^k_{e}\geq 0,~ \mbox{ for } e \in \mathcal{A},~ k=1,2,\ldots,K\\
\end{align*}

Notes:

K=1 is ordinary single-commodity network flow. Integer solutions for free when node-supplies and arc capacities are integer.
K=2 example below with integer data gives a fractional basic optimum. This example doesn't have any feasible integer flow at all.

\begin{remark}{}{}
Unfortunately, the same integrality theorem does not hold in the multi-commodity network flow probem.  Nontheless, if the quanties in each flow are very large, then the LP solution will likely be very close to an integer valued solution.
\end{remark}
